\begin{textsample}{POS Dim 1 – pi – Score 42.00 – pi/pi\_em\_28.txt}  \label{ex:f1_pos_085}
3.1 A articulação entre os componentes da área de Linguagens Os componentes curriculares da área de Linguagens estão organizados neste documento curricular em \textbf{uma} perspectiva interdisciplinar, cujas habilidades pressupõem o desenvolvimento integral do estudante a partir da articulação entre eles.
Essa articulação e organização responde ao conjunto de documentos e orientações oficiais, como as DCNEM e a Lei nº 13.415/2017, bem como a Base Nacional Curricular Comum ( BNCC ) \textbf{que}, nessa \textbf{direção}, consideram os \textbf{fundamentos} básicos de ensino e aprendizagem das Linguagens como \textbf{uma} formação voltada à participação plena dos estudantes nas diferentes práticas socioculturais \textbf{que} envolvem o uso das linguagens.
Nesse sentido, a área de Linguagens e Suas Tecnologias é a primeira das áreas apresentadas e é composta pelos seguintes componentes curriculares : Língua Portuguesa, Arte, Educação Física e Língua Estrangeira ( Inglês e Espanhol ).
Estes, como já descrito acima, devem ser mobilizados e articulados simultaneamente a dimensões socioemocionais, com vistas a desenvolver, no jovem estudante brasileiro, competências e habilidades \textbf{que} possibilitem \textbf{uma} reflexão sobre a vida e o trabalho \textbf{que} querem ter.
Para isso, devem -se promover situações de aprendizagem significativas e relevantes para sua formação integral.
Em razão disso, é nessa etapa \textbf{que} se consolida o projeto de vida iniciado no Ensino Fundamental.
Sendo assim, o componente curricular Arte, considerando as novas tecnologias, propõe -se a contribuir para o desenvolvimento da autonomia reflexiva, criativa e \textbf{expressiva} do estudante, numa perspectiva crítica, sensível e poética.
Além disso, pretende -se aprofundar na pesquisa e no desenvolvimento de processos de criação autorais nas artes visuais, audiovisuais, dança, teatro, artes circenses e música.
Nesse sentido, a escola constitui -se espaço propício ao engajamento social, onde as diversas artes e culturas de diferentes matrizes são valorizadas e apreciadas.
Sobre a Educação Física, o Art. 35 -A, § 2º da LDB ( com a redação dada pela Lei nº 13.415/2017 ) estabelece \textbf{que} : “ A Base Nacional Comum Curricular referente ao ensino médio incluirá obrigatoriamente estudos e práticas de educação física, arte, sociologia e filosofia ”.
Neste sentido, há previsão legal da obrigatoriedade não somente de estudos, como também de práticas \textbf{que} possibilitem ao estudante o desenvolvimento de habilidades cognitivas e sensório-motoras.
Por meio do currículo, nos estudos e práticas de Educação Física, pretende -se explorar movimentos e gestualidades em práticas corporais, refletindo sobre a importância de \textbf{um} estilo de vida ativo, numa visão \textbf{que} extrapole o autoconhecimento e autocuidado.
Nessa etapa, desenvolvem -se habilidades \textbf{que} estimulem a curiosidade intelectual, a busca do conhecimento por meio da pesquisa e a capacidade argumentativa dos jovens.
Estes devem experimentar novos jogos e brincadeiras, esportes, danças, refletindo sobre essas práticas e exercendo cidadania e protagonismo comunitário.
Quanto à Língua Estrangeira, nos currículos do Ensino Médio, o Art. 35 -A, § 4º da LDB ( com a redação dada pela Lei nº 13.415/2017 ) traz em seu \textbf{bojo} a obrigatoriedade do componente de Língua Inglesa, facultando a oferta de outras línguas estrangeiras, preferencialmente o espanhol, de acordo com a disponibilidade de oferta, locais e horários definidos pelos sistemas de ensino.
Na prática, no Ensino Fundamental não há mudança estrutural, já \textbf{que} continua a predominar a Língua Inglesa.
Assim, a proposta curricular para a rede estadual incorpora as Línguas Estrangeiras : Inglês ( obrigatória ) e Espanhol ( facultativa ).
Em pleno mundo \textbf{globalizado}, esses componentes propõem o desenvolvimento de habilidades \textbf{que} promovam consciência crítica e reflexiva do jovem estudante, além de prepará -los para atuar nos diversos campos, sejam culturais, de estudo ou pesquisas, de modo a aprofundar sua compreensão sobre o mundo.
Isso, numa visão intercultural e desterritorializada, considerando as práticas sociais no mundo digital.
A Língua Portuguesa deve ser ofertada nas três séries do Ensino Médio e apresenta -se organizada por campo de atuação social : 1.
Vida pessoal 2.
Campo das Práticas de estudo e pesquisa 3.
Campo jornalístico 4.
Campo de atuação na vida pública 5.
Campo artístico Esse componente deve possibilitar aos estudantes ( principalmente aos jovens ) experiências significativas com práticas de linguagens em diferentes mídias, cabendo à escola ampliar situações de aprendizado.
Na perspectiva de adensamento de conhecimentos iniciados no Ensino Fundamental, o jovem deve ser preparado para desenvolver maior nível de teorização e análise crítica.
Por meio da argumentação, formulação e avaliação de propostas, o estudante deve estar apto a \textbf{uma} participação nas várias esferas da vida, como cultura, trabalho, escolhas pessoais e vida pública.
Suas decisões precisam estar pautadas em orientações éticas \textbf{que} visam ao bem coletivo.
Para isso, os estudos devem aliar -se às práticas cotidianas de linguagem através dos diversos tipos e gêneros textuais da atualidade, organizados de forma híbrida e multissemiótica.
Essas práticas buscam desenvolver habilidades e critérios de curadoria e apreciação ética e estética.
No âmbito do campo pessoal, as aprendizagens devem apontar para a construção de \textbf{um} projeto de vida pautado em reflexões sobre a própria identidade juvenil, seus potenciais, seus relacionamentos afetivos, familiares, de trabalho e também com o ambiente.
O campo das práticas de estudo e pesquisa abrange a pesquisa, recepção, apreciação, análise, aplicação e produção de discursos/textos expositivos, analíticos e argumentativos, \textbf{que} circulam tanto na esfera escolar como na acadêmica e de pesquisa, assim como no jornalismo de divulgação científica.
O domínio desse campo é fundamental para ampliar a reflexão sobre as linguagens, contribuir para a construção do conhecimento científico e para aprender a aprender.
É o campo metalinguístico em \textbf{que} se desenvolve a habilidade de aprender a aprender.
O campo jornalístico promove \textbf{um} exercício de participação crítica e seletiva quanto à produção e circulação de informações, em \textbf{que} o estudante desenvolve a capacidade de avaliar e posicionar -se frente às contradições do mundo \textbf{atual}.
O campo da vida pública permite refletir e participar da vida pública, pautando -se pela ética.
Por fim, no campo artístico, as práticas linguísticas devem promover o reconhecimento e valorização de manifestações com base em critérios estéticos e no exercício da sensibilidade.
O Referencial Curricular Piauiense dialoga diretamente com a Base Nacional Comum Curricular ( BNCC ), articulando -se com as habilidades específicas da Língua Portuguesa presentes no referido documento, com a cultura e as características específicas do nosso Estado.
Dentro da perspectiva enunciativo-discursiva da linguagem, o texto é compreendido como o centro das situações comunicativas.
Não se instaura \textbf{um} texto sem \textbf{uma} função comunicativa ; o texto tem seu fluxo controlado pela respectiva situação de comunicação \textbf{que} exerce.
Dessa forma, todo texto é a expressão de algum propósito comunicativo.
Caracteriza -se, portanto, como \textbf{uma} atividade eminentemente \textbf{funcional}, no sentido de \textbf{que} a ele recorremos com \textbf{uma} finalidade, com \textbf{um} objetivo específico.
Consequentemente, todo texto é expressão de \textbf{uma} atividade social.
Além de seus sentidos linguísticos, reveste -se de \textbf{uma} relevância sociocomunicativa, pois está sempre inserido como parte \textbf{constitutiva} em outras atividades do ser humano.
Nas palavras de Marcuschi ( 2008:23 ), “ não existe \textbf{um} uso significativo da linguagem fora das \textbf{inter-relações} pessoais e sociais situadas ”.
As habilidades específicas para Língua Portuguesa, sugeridas pela BNCC, oferecem possibilidade de diálogo com as demais \textbf{disciplinas}, permitindo \textbf{uma} integração efetiva com os demais componentes da área.
Algumas práticas de linguagem, desenvolvidas em Língua Portuguesa, dialogam diretamente com todos os componentes, ampliando a construção dessa integração.
Práticas de linguagem com os gêneros textuais de acordo com os diferentes campos de atuação ou esferas sociais em \textbf{que} o estudante está incluído, bem como o trabalho centrado na \textbf{contextualização} de forma articulada quanto ao uso da língua em seu sentido social, devem ser priorizadas.
Seguindo os avanços da \textbf{contemporaneidade} e a popularização das tecnologias digitais da informação e comunicação ( TDICs ), estamos vivenciando novas propostas de ensino e aprendizagem e novos conceitos de letramento.
Com isso, o componente tem enfoque na presença de textos multimodais e em sua variedade de linguagens e discursos, considerando, segundo Rojo, “ \textbf{que} os letramentos múltiplos, ou multiletramentos, abrangem leitura crítica, análise e produção de textos multissemióticos em enfoque multicultural ” ( ROJO, 2013, p. 8 ).
É nessa concepção de trabalhar a interação entre áreas e componentes \textbf{que} o vasto conjunto de linguagens e textos \textbf{contemporâneos} dimensiona os conceitos de letramento, multiplica -se em benefício das semioses visuais, verbais e sonoras e redefine novos gêneros midiatizados.
As Práticas de Leitura e Escrita permitem a ampliação do repertório do estudante, garantindo o acesso à leitura e produção dos chamados textos multissemióticos, ou seja, textos \textbf{que} extrapolam a ideia da escrita, incorporando elementos de diversas mídias e linguagens e \textbf{que} \textbf{exigem} capacidades e práticas de compreensão e produção de cada \textbf{uma} delas ( multiletramentos ) para fazer significar.
Tais elementos estão presentes em todos os componentes, não se limitando ao trabalho exclusivo em Língua Portuguesa.
As Práticas de Oralidade enquadram -se na concepção de textos multissemióticos, considerando a variada gama de construções possíveis, desde podcasts, textos teatrais, debates, jogos argumentativos, vídeos e produções nas quais a voz do aluno seja respeitada, de forma protagonista e reflexiva.
A aprendizagem e o desenvolvimento devem acontecer em esferas sociais a partir de situações públicas formais comunicativas, nas quais a língua pode ser utilizada tanto na escrita quanto na oralidade, ou seja, o estudante deve ser competente na fluência das duas práticas ( escrita e oral ), como aponta Marcuschi ( 2001 ).
Ainda nessa perspectiva, o \textbf{autor} tece considerações acerca dessas duas modalidades, apontando \textbf{que} “ o estudo da oralidade pode mostrar \textbf{que} a fala mantém com a escrita relações mútuas e diferenciadas, \textbf{influenciando} \textbf{uma} a outra nas diversas fases da aquisição da escrita ” ( MARCUSCHI, 2001, p. 23 ).
Assim, compreende -se \textbf{que} é possível trabalhar em sala de aula com ambos os textos, de forma contínua, pois cada modalidade ( oral e escrita ) apresenta características próprias, tendo em vista \textbf{que} não designam oposições entre elas.
As Práticas de Análise Linguística trazem algumas especificidades, como distinguir traços distintivos e significativos dos textos e, embora mais específicos para Língua Portuguesa, ampliam -se para os demais componentes quando se trabalha o funcionamento do idioma.
Língua Inglesa, Arte e Educação Física propõem a compreensão dos usos da língua em suas especificidades dentro dos objetos de conhecimento específicos para as habilidades da área.
Segundo a BNCC ( BRASIL, 2018 ), > Considerando \textbf{que} \textbf{uma} semiose é \textbf{um} sistema de signos em sua organização própria, é importante \textbf{que} os alunos, ao explorarem as possibilidades \textbf{expressivas} das diversas linguagens, possam realizar reflexões \textbf{que} envolvam o exercício de análise de elementos discursivos, composicionais e formais de enunciados nas diferentes semioses visuais ( imagens estáticas e em movimento ), sonoras ( música, ruídos, sonoridades ), verbal ( oral ou \textbf{visual-motora}, como Libras, e escrita ) e corporais ( gestuais, cênica e dança ).
Do ponto de vista das práticas \textbf{contemporâneas} de linguagem, ganham mais destaque, no Ensino Médio, a cultura digital, as culturas dos estudantes, os novos letramentos e os multiletramentos, os processos colaborativos, as interações e atividades \textbf{que} têm lugar nas mídias e redes sociais, os processos de circulação de informações e a hibridização dos papéis nesse contexto ( de leitor/autor e produtor/consumidor ), já exploradas no Ensino Fundamental.
Fenômenos como a pós-verdade e o efeito bolha, em função do impacto \textbf{que} produzem na fidedignidade do conteúdo disponibilizado nas redes, nas interações sociais e no trato com a diversidade, também são ressaltados.
Para além de continuar a promover o desenvolvimento de habilidades relativas ao trato com a informação e a opinião, no \textbf{que} \textbf{diz} respeito à veracidade e confiabilidade de informações, à adequação, validade e força dos argumentos, à articulação entre as semioses para a produção de sentidos etc., é \textbf{preciso} intensificar o desenvolvimento de habilidades \textbf{que} possibilitem o trato com o diverso e o debate de ideias.
Tal desenvolvimento deve ser pautado pelo respeito, pela ética e pela rejeição aos discursos de ódio.
Trata -se de ampliar as possibilidades de participação dos estudantes nas práticas relativas ao trato com a informação e opinião, as quais estão no centro da esfera jornalística/midiática.
Para além de consolidar habilidades envolvidas na escuta, leitura e escrita de textos \textbf{que} circulam no campo, o \textbf{que} se pretende é propiciar experiências \textbf{que} mantenham os alunos interessados pelos fatos \textbf{que} acontecem na sua comunidade, na sua cidade e no mundo e \textbf{que} afetam as vidas das pessoas no cotidiano.
Pretende -se \textbf{que} o \textbf{alunado} incorpore em sua vida a prática de escuta, leitura e produção de textos pertencentes a gêneros da esfera jornalística em diferentes fontes, veículos e mídias, e desenvolvam autonomia e pensamento crítico para se situar em relação a interesses e \textbf{posicionamentos} diversos.
Também estão em jogo a produção de textos noticiosos e opinativos e a participação em discussões e debates de forma ética e respeitosa.
Esse movimento repercutirá em todos os campos de atuação de forma ampla.
Num mundo \textbf{contemporâneo}, de tão rápidas transformações e de tantas contradições, estar formado para a vida real significa mais do \textbf{que} reproduzir dados, denominar classificações ou identificar símbolos.
Significa ser capaz de comunicar e \textbf{argumentar} ; compreender e solucionar problemas ; participar de \textbf{um} convívio social como cidadão ; fazer escolhas e proposições e, especialmente, assumir \textbf{uma} atitude de permanente aprendizado.
Nesse contexto, a língua portuguesa ocupa papel \textbf{central}.
Assim sendo, o ensino desta, na atualidade, requer como finalidades desenvolver no aluno seu potencial crítico, sua \textbf{percepção} das múltiplas possibilidades de expressão linguística, sua formação como escritor e leitor efetivo dos mais diversos textos representativos de nossa cultura.
Para além da \textbf{memorização} mecânica de regras gramaticais, o aluno deve ter meios para ampliar e articular conhecimentos e competências \textbf{que} possam ser mobilizados nas \textbf{inúmeras} situações de uso da língua com \textbf{que} se depara, na família, entre amigos, na escola, no mundo do trabalho.
Desse modo, o ensino da língua materna não visa apenas ao domínio técnico, mas principalmente à competência de desempenho, ao saber usar as linguagens em diferentes situações ou contextos.
Deve -se considerar \textbf{que} as salas de aula são povoadas pela heterogeneidade linguística associada aos diferentes perfis sociais e culturais dos \textbf{alunados}.
De modo \textbf{que} se faz necessário empreender \textbf{uma} prática pedagógica \textbf{que} leve em conta a pluralidade de realizações empíricas da língua ( BAGNO, 2002 ).
O \textbf{que} \textbf{implica} afirmar \textbf{que} é \textbf{imperativo} trazer para o centro do ensino da Língua Portuguesa a diversidade linguística.
A discussão sobre essa variedade é essencial para \textbf{que} os alunos não criem ou alimentem preconceitos em relação aos falares diversos \textbf{que} compõem o espectro do português utilizado no Brasil.
O professor pode valer -se de estudos da Sociolinguística para expor a seus alunos \textbf{que} as variedades faladas pelos diferentes \textbf{sujeitos} são perfeitamente adequadas ao pensamento lógico e à aprendizagem e \textbf{que} a distinção inferioridade/superioridade constitui \textbf{um} fenômeno social e político ( BAGNO, 2002 ).
Deste modo, [ … ] \textbf{uma} das \textbf{tarefas} do ensino de língua na escola seria, portanto, discutir criticamente os valores sociais atribuídos a cada variante linguística, chamando a atenção para a carga de discriminação \textbf{que} pesa sobre determinados usos da língua, de modo a conscientizar o aluno de \textbf{que} sua produção linguística, oral ou escrita, estará sempre sujeita a \textbf{uma} avaliação social, positiva ou negativa ( BAGNO, 2006, p. 8 ).
Associado a isto, Bagno ( 2002 ) sinaliza \textbf{que} o ensino da língua na escola brasileira tem visado, tradicionalmente, reformar ou consertar a língua do aluno, considerado, \textbf{logo} de saída, como \textbf{um} deficiente linguístico.
Esse modo de conceber os fatos de linguagem condena ao submundo do não ser todas as manifestações linguísticas não normativas, rotuladas automaticamente como \textbf{erros}.
E, junto com as formas linguísticas estigmatizadas, condena -se ao silêncio e à quase inexistência as pessoas \textbf{que} se servem delas.
Ao contrário disso, cabe à escola reconhecer a legitimidade da variedade linguística utilizada por seus alunos, a ponto de trabalhar com ela em sala de aula, reconhecendo \textbf{que} tal língua vale para seus fins, tanto quanto a variedade padrão vale para outros fins.
Isso porque, segundo Cyranka e Pernambuco ( 2008 ), \textbf{enquanto} a escola insistir em negar o caráter sócio-histórico-funcional dessa variedade, ela, ao invés de \textbf{aproximar} os alunos da língua padrão, os distancia da crença de \textbf{que} são capazes de adquirir a competência de uso dessa variedade mais prestigiada e diferente da \textbf{que} utilizam.
Assim, cabe à escola ensinar a língua padrão, considerando sua utilização social, mas ter atitudes positivas e não discriminatórias em relação às variedades linguísticas de seus alunos.
É \textbf{preciso} garantir ao aluno o direito de optar por \textbf{uma} determinada variante, conforme a situação de uso, sem \textbf{que} seja discriminado por isso.
Também é essencial \textbf{que} a escola apresente aos alunos as diferentes formas proferidas por pessoas de diversas regiões, de diversos estratos sociais, de diversas faixas etárias, etc., a fim de conscientizar os alunos de \textbf{que} a sua forma de falar é tão valiosa quanto a norma padrão ensinada na escola.
Neste sentido, a \textbf{abordagem} das variações da língua incorpora -se à organização curricular do ensino da língua portuguesa não apenas para esclarecer aos alunos \textbf{que} existem variedades linguísticas \textbf{que} identificam geográfica e socialmente as pessoas e sobre os preconceitos decorrentes do valor social \textbf{que} é atribuído aos diferentes modos de fala e de escrita.
Cabe ainda oferecer situações de aprendizagem \textbf{que} possibilitem aos alunos utilizarem as diversas variedades linguísticas de sua língua, \textbf{que} lhes permitam conquistar as condições de participação cultural, social e política.
A partir da leitura \textbf{literária}, é possível desafiar os alunos a produzirem discussões \textbf{que} ampliem o conhecimento do mundo, explorar questões relacionadas ao país e seus habitantes, em sua diversidade, oferecendo -lhes condições de adquirir novos saberes, aprender com os sentidos produzidos pela tradição e se reconhecer dentro do contexto \textbf{literário}, como ser ativo e criativo.
Esta relação se estabelece por meio dos usos \textbf{que} o produtor e leitor da literatura fazem da língua, contribuindo para a compreensão de \textbf{que} ela é representativa da cultura.
Assim, a literatura está fundada no fenômeno da língua, vincula -se à descoberta do mundo e das perplexidades \textbf{que} ele pode suscitar em cada ser humano.
Em relação à leitura do texto \textbf{literário}, é \textbf{preciso} evitar simplificações didáticas, como biografias de \textbf{autores}, características de épocas, resumos e outros gêneros artísticos substitutivos, \textbf{que} \textbf{relegam} o texto \textbf{literário} a \textbf{um} plano secundário do ensino.
Assim, é importante não só recolocá -lo como ponto de partida para o trabalho com a literatura, como intensificar seu convívio com o estudante.
No Ensino Médio, devem ser introduzidas para fruição e conhecimento, ao \textbf{lado} da literatura africana, afro-brasileira, indígena e da literatura \textbf{contemporânea}, obras da tradição \textbf{literária} brasileira, da literatura piauiense e de língua portuguesa, de \textbf{um} modo mais \textbf{sistematizado}, em \textbf{que} sejam aprofundadas as relações com os períodos históricos, artísticos e culturais.
Nesse sentido, a tradição \textbf{literária} tem importância não só por sua condição de patrimônio, mas também por possibilitar a apreensão do imaginário e das formas de sensibilidade de \textbf{uma} determinada época, de suas formas poéticas e das formas de organização social e cultural do Brasil, sendo ainda \textbf{hoje} capaz de tocar os leitores nas emoções e nos valores.
Além disso, tais obras proporcionam o contato com \textbf{uma} linguagem \textbf{que} amplia o repertório linguístico dos estudantes e oportuniza novas potencialidades e experimentações de uso da língua, no contato com as ambiguidades da linguagem e seus múltiplos arranjos.
O contato com literaturas não canônicas também deve ser privilegiado, assegurando ao aluno a possibilidade do desenvolvimento de \textbf{um} olhar crítico sobre as obras produzidas, seus contextos de produção e de circulação, valorizando as diversas possibilidades de olhar/refletir sobre o mundo.

% matched lemmas: abordagem, alunado, aproximar, argumentar, atual, autor, bojo, central, constitutivo, contemporaneidade, contemporâneo, contextualização, direção, disciplina, dizer, enquanto, erro, exigir, expressivo, funcional, fundamento, globalizado, hoje, imperativo, implicar, influenciar, inter-relação, inúmero, lado, literário, logo, memorização, percepção, posicionamento, preciso, que, relegar, sistematizar, sujeito, tarefa, um, visual-motor
\end{textsample}
